\begin{minipage}[t][][t]{0.70\linewidth}
    À fin de vérifier la valeur de la concentration molaire $C$ de la solution
    $\mathrm{(S)}$, on réalise le dosage d'un volume $V_{A}=\qty{10}{\mL}$ de la
    solution $\mathrm{(S)}$ par une solution aqueuse $\mathrm{(S_{B})}$
    d'hydroxyde de sodium $\ce{Na^+_{(aq)} + HO^-_{(aq)}}$ de concentration
    molaire ${C_{B}=\qty{1e-2}{\mol\per\L}}$ en utilisant le montage de la
    Fig.~\ref{fig:2bac_svt_national_2024_rattrapage_chim_2}.  La courbe de la
    Fig.~\ref{fig:2bac_svt_national_2024_rattrapage_chim_3} représente les
    variations du $\mathrm{pH}$ du mélange en fonction du volume $V_{B}$ de la
    solution $\mathrm{(S_{B})}$ ajoutée.

    \begin{enumerate}
        \item Attribuer au numéros~(1),~(2) et~(3) du montage de la
              Fig.~\ref{fig:2bac_svt_national_2024_rattrapage_chim_2} les noms
              correspondants.~\textbf{(0,75pt)}
        \item Écrire l'équation de la réaction du dosage supposée
              totale.~\textbf{(0,5pt)}
        \item Déterminer les coordonnées ${(V_{B\mathrm{E}}~;~\mathrm{pH_{E}})}$
              du point d'équivalence.~\textbf{(1pt)}
        \item Calculer la concentration $C$ de la solution
              $\mathrm{(S)}$.~\textbf{(1pt)}
        \item Préciser, parmi les indicateurs colorés ci-dessous,
              l'indicateur convenable pour repérer cette équivalence.
              Justifier.~\textbf{(0,5pt)}
    
              \begin{minipage}{\linewidth}
                  \begin{table}[H]
                      \centering
                      \renewcommand{\arraystretch}{1.1}
                      \begin{tabularx}{\textwidth}{|p{4cm}|C|C|C|}
                          \hline
                          \textbf{Indicateur coloré} &
                          \textbf{Teinte acide} &
                          \textbf{Zone de virage} &
                          \textbf{Teinte basique} \\
                          \hline
                          Rouge de méthyle & Rouge & $4,4-6,2$ & Jaune \\
                          \hline
                          Rouge de crésol & Jaune & $7,2-8,8$ & Rouge \\
                          \hline
                          Alizarine & Rouge & $11,0-12,4$ & Violette \\
                          \hline
                      \end{tabularx}
                  \end{table}
              \end{minipage}
    \end{enumerate}
    
\end{minipage}
\hfill
\begin{minipage}[t][][t]{0.28\linewidth}
    \begin{figure}[H]
        \centering
        \includegraphics[width=\textwidth]{2bac_svt_national_2024_rattrapage_chim_2}
        \caption{}
        \label{fig:2bac_svt_national_2024_rattrapage_chim_2}
    \end{figure}
    \begin{figure}[H]
        \centering
        \includegraphics[width=\textwidth]{2bac_svt_national_2024_rattrapage_chim_3}
        \caption{}
        \label{fig:2bac_svt_national_2024_rattrapage_chim_3}
    \end{figure}
\end{minipage}
